\documentclass[aspectratio=169]{beamer}
%[aspectratio=169]   \usepackage[czech]{babel}
\usepackage{apo-lecture-en}
\usepackage{pdfpages}
\usepackage{pdfcomment}
\usepackage{listings}
\usepackage{array,multirow}

\title{Virtual memory in RISC-V}
\subtitle{LinuxDays 2026}
\author{Pavel Píša \phantom{xxxxxxxxx} Petr Štěpán \\ \small\texttt{pisa@fel.cvut.cz}\phantom{xxxx}\small\texttt{stepan@fel.cvut.cz}\\ \phantom{xxx} \\
based on B35APO slides \\
and the bachelor thesis virtual memory for QtRvSim by Adam Němec  \\
License: CC-BY-SA}
\begin{document}

\maketitle

\section{QtRvSim}

\begin{frame}
\frametitle{QtRVSim -- RISC-V CPU Simulator for Education Purposes}

\begin{columns}
\begin{column}{0.68\textwidth}

\begin{itemize}
\item RISC-V simulator with processor pipeline, cache and branch predictor visualization
\item An tool to support understanding computer architectures principles
\item Authors: Karel Kočí, Pavel Píša, Jakub Dupák, Jiří Štefan, Adam Němec
\item It will be used for interactive demonstrations during presentation
\item Online simulator version, links to downloads and related \\ Czech Technical University in Prague materials and courses links \\ \url{https://comparch.edu.cvut.cz/}
\end{itemize}

\end{column}
\hfill
\begin{column}{0.30\textwidth}

\begin{center}
\includegraphics[width=0.9\linewidth]{generated/simulator/qtrvsim-logo.pdf}
\end{center}

\end{column}
\hfill
\end{columns}
\end{frame}


\begin{frame}
\frametitle{QtRVSim - Download and Presentations}
\begin{itemize}
\item Windows, Linux, Mac \url{https://github.com/cvut/qtrvsim/releases}
\item  Ubuntu \url{https://launchpad.net/~qtrvsimteam/+archive/ubuntu/ppa}
\item Suse, Fedora and Debian \url{https://software.opensuse.org/download.html?project=home\%3Ajdupak\&package=qtrvsim}
\item Online version \url{https://comparch.edu.cvut.cz/qtrvsim/app/}
\item \textit{QtRvSim - RISC-V Simulator with Cache and Pipeline Visualization}, RISC-V International Academic and Training SIG meeting, 2023, recording \url{https://youtu.be/J6AcPZZ_ISg}
\item \textit{QtRVSim—Education from Assembly to Pipeline, Cache Performance, and C Level Programming}, FOSDEM 2023, Brussels, \url{https://fosdem.org/2023/schedule/event/rv_qtrvsim/}
\item Dupák, J.; Píša, P.; Štepanovský, M.; Kočí, K.: \textit{QtRVSim – RISC-V Simulator for Computer Architectures Classes}, In: embedded world Conference 2022, Haar: WEKA FACHMEDIEN GmbH, 2022. p. 775-778. ISBN 978-3-645-50194-1. (\href{https://comparch.edu.cvut.cz/publications/ewC2022-Dupak-Pisa-Stepanovsky-QtRvSim.pdf}{PDF}) (\href{https://comparch.edu.cvut.cz/slides/ewc22-qtrvsim.pdf}{Slides})

\end{itemize}
\end{frame}

\section{Processor}


\begin{frame}
\frametitle{Základy ochrany programů}

\begin{itemize}
\item Potřebujeme zabezpečit program a operační systém tak, aby chyba uživatelských programů nemohla narušit běh ostatních programů.
\begin{itemize}
\item Je nutné zajistit bezpečnost dat - nemožnost přečíst tajné informace, klíče, atd.
\item Všechna data musí být dosažitelná pouze oprávněným uživatelem.
\end{itemize}
\item Procesor potřebuje zakázat běžným programům měnit nastavení systému, které je klíčové pro vzájemnou ochranu uživatelských procesů a JOS.
\begin{itemize}
\item v příští lekci se podíváme co je stránkování a jak lze pomocí stránkování zabezpečit přístup k paměti
\end{itemize}
\item Různé procesory používají k ochraně programů různé přístupy.
\item V této lekci si probereme přístup RISC V procesoru
\end{itemize}
\end{frame}


\begin{frame}
\frametitle{Módy procesoru RISC V}

RISC V procesor má tři úrovně ochrany (o virtualizaci nemluvě):
\begin{itemize}
\item Machine Mode (M-mode): Nejvyšší úroveň oprávnění. V tomto režimu běží kód s přímým přístupem k hardwaru (bootloader, firmware, zabezpečené monitorovací programy). M-mode je povinný pro všechny implementace RISC-V.
\item  Supervisor Mode (S-mode): Střední úroveň oprávnění, určená pro běh operačních systémů. S-mode může využívat virtuální paměť a mechanismy ochrany paměti.
\item User Mode (U-mode): Nejnižší úroveň oprávnění, ve které běží uživatelské aplikace. Aplikace v tomto režimu nemají přímý přístup k hardwaru a pro privilegované operace musí volat služby operačního systému (v S-mode)
\end{itemize}

Všechny procesory RISC V musí implementovat M-mode, ostatní úrovně nemusí být implementovány.
\end{frame}


\begin{frame}
\frametitle{Změny módu procesoru RISC V}

Jak přechází RISC V mezi úrovněmi ochrany:
\begin{itemize}
\item Zvýšení úrovně
\begin{itemize}
\item Přechod U $\rightarrow$ S, U $\rightarrow$ M nebo S $\rightarrow$ M
\item Nejnebezpečnější jsou přechody U $\rightarrow$ S, U $\rightarrow$ M.
\begin{itemize}
\item Musíme zabránit běžet uživatelskému programu v módu S nebo M
\end{itemize}
\item Kdy je potřeba použít S nebo M mód?
\begin{itemize}
\item Když potřebujeme vyřešit spolupráci s periférií na základě přerušení.
\item Když uživatelský program chce použít HW (zapsat/číst do/z souboru, poslat paket po síti, atd.) --
Pro tuto činnost potřebujeme speciální instrukci ecal - systémové volání.
\end{itemize}
\item Nesmíme dovolit spustit uživatelský program v módu M nebo S
\item Adresa pro přechod do S módu je uvedena v registru stvec a pro přechod do M módu mtvec
\item lze nastavit, zda se vybrané přerušení bude obsluhovat v M nebo S módu
\end{itemize}
\end{itemize}

\end{frame}



\begin{frame}
\frametitle{Módy procesoru RISC V}

Jak přechází RISC V mezi úrovněmi ochrany:
\begin{itemize}
\item Snížení úrovně
\begin{itemize}
\item Přechod M $\rightarrow$ S, M $\rightarrow$ U nebo S $\rightarrow$ U
\item Toto není tak nebezpečně, snížení úrovně je vždy lepší.
\begin{itemize}
\item Snížení úrovně znamená omezení práv
\end{itemize}
\item Speciální instrukce mret a sret, které zajistí přechod do nižší úrovně.
\item Úroveň je nastavena v systémovém registru sstatus.mpp nebo sstatus.spp.
\item Adresa, kam se má skočit instrukcí je v registru mepc nebo sepc.
\item Tento přechod je vlastně návratem s obsluhy přerušení, nebo systémového volání
\begin{itemize}
\item I v případě startu, kdy nikdo systémové volání nezavolal.
\end{itemize}
\end{itemize}
\end{itemize}

\end{frame}


\begin{frame}
\frametitle{Přepnuí do U módu}

\begin{itemize}
 \item Po zapnutí procesoru je vždy nastaven M mód (Machine mód)
\item Přepnutí do U módu se provede takto
 \begin{itemize}
 \item do registru mstatus nastavíme, že se chceme přepnout do U módu, hodnota 0 v poli MPP (Machine Previous Privilige)
\begin{itemize}
 \item Poznámka: register slouží k zaznamenání módu, ze kterého jsme se dostali do M módu, abychom se tam mohli vrátit
 \item Protože náš procesor se právě spustil, přepnutí do U módu simulujeme jakoby návratem do U módu instrukcí mret
 \item Pokud bychom se chtěli přepnout do S módu nastavili bychom do MPP 1
\end{itemize}
 \item do registru mepc nastavíme, jakou adresu chceme v U módu spustit
\begin{itemize}
 \item opět tento register slouží k zaznamenání adresy, ze které například přerušením se spustil program v M módu
\end{itemize}
 \item provedení instrukce mret nás přepne do U módu na adresu mepc
\end{itemize}
\end{itemize}
\end{frame}


\begin{frame}
\frametitle{Přepnutí do S/M módu příklad}

Pro návrat do S nebo M módu z U módu se musí nejdříve v M/S módu nastavit
\begin{itemize}
\item stvec nebo mtvec adresu programu, který se bude spouštět při všech přepnutích
 \begin{itemize}
 \item jedná se o jeden program, který zjistí co se stalo (HW přerušení, vyjímka, ecall) a v případě ecall podle hodnoty registru a7 rozlišuje, co uživatel chce provést
 \end{itemize}
\item pokud chceme přecházet do S módu, musíme nastavit v registru medeleg přepnutí obsluhy v S módu
\item na závěr povolíme přechody do S nebo M módu v registru sstatus, nebo mstatus
\end{itemize}

Uživatelský program:
\begin{itemize}
\item provedení instrukce ecall přepne procesor do S/M módu a spustí nastavený program
\end{itemize}
\end{frame}


\begin{frame}
\frametitle{Základy ochrany OS}

\begin{itemize}
\item Uživatelský proces nemůže přímo přistupovat k HW počítače.
\item Uživatelský proces musí požádat OS o zpřístupnění HW, nebo vyřízení požadavku.
\item To co OS umožňuje procesům jsou pouze systémová volání (system calls).
\item Na rozdíl od volání funkce neznáme adresy funkcí v jádře.
\item Systémové volání se vybírá číslem systémového volání.
\item Vlastní zavolání funkce je vyvolání přerušení nebo výjimky specializovanou instrukcí:
\begin{itemize}
\item x86 -- používá přímo int 0x80 -- vyvolej přerušení číslo 0x80.
\item x86 -- nověji specializované instrukce sysenter/syscall -- podobný mechanizmus, ale nevyužívá uživatelský zásobník, složitý mechanizmus bran a přistupuje k méně paměti $\rightarrow$ je rychlejší.
\item RISC-V -- specializovaná instrukce ecall -- obdobně jako při přerušení.
\end{itemize}
\item Při přerušení/výjimce se volá obslužná rutina v privilegovaném řežimu (typicky systémovém) z adresy, která je nastavená OS (na RISC-V registry \texttt{mtvec}/\texttt{stvec}), uživatel ji nemůže měnit.
\end{itemize}
\end{frame}


\begin{frame}
\frametitle{Systémové volání}

Uživatelský program nastaví parametry volání do registrů vyvolá speciální přerušení/výjimku (na RISC-V \texttt{ecall}).

OS podle čísla systémového volání zavolá odpovídající funkci.

\begin{center}
  \includegraphics[width=0.95\textwidth]{generated/io/irq/irq-and-os-cz.pdf}
\end{center}

Návrat ze systémového volání odpovídá návratu z přerušení/výjimky (na RISC-V většinou \texttt{sret}).
\end{frame}


\begin{frame}[shrink=5]
\frametitle{ABI systémového volání}

\begin{itemize}
\item API -- application programming interface -- popis, jaké funkce může Váš program zavolat z knihoven.
\begin{itemize}
\item API je definováno pro jazyk C hlavičkovými soubory.
\item API definuje i co dané funkce dělají, jaké vrací hodnoty, jak se chovají při chybě
\item zkuste např. \texttt{man 2 read}
\end{itemize}
\item ABI -- application binary interface -- popis jaké registry a jaké instrukce použít
\item ABI pro systémová volání na architektuře RISC-V
\begin{itemize}
\item \texttt{a7} -- obsahuje číslo systémového volání (přehled např. \url{https://jborza.com/post/2021-05-11-riscv-linux-syscalls/} nebo \url{https://marcin.juszkiewicz.com.pl/download/tables/syscalls.html})
\item \texttt{a0} až \texttt{a5} -- parametry systémového volání (Linuxové systémové volání má maximálně 6 parametrů)
\item systémové volání se provede instrukcí \texttt{ecall}
\item \texttt{a0} obsahuje návratovou hodnotu systémového volání
\begin{itemize}
\item pokud by mělo vracet více dat (např. čtení ze souboru) musí uživatel zadat ukazetel na buffer a velikost bufferu, kam OS zapíše data.
\end{itemize}
\end{itemize}
\end{itemize}
\end{frame}

\begin{frame}[fragile]
\frametitle{Příklad Hello world!}

\begin{minted}[fontsize=\footnotesize]{gas}
.global _start
.text
_start:
#   write(1, "Hello world!\n", 13);
   addi  a7, zero, 64
   addi  a0, zero, 1
   la    a1, zero, text_1
   addi  a2, zero, 13
   ecall
final:
#   exit(0);
   addi  a7, zero, 93
   addi  a0, zero, 0
   ecall
   ebreak
   j final
.data
# store ASCII text, no termination
text_1: .ascii "Hello world!\n"
\end{minted}
\end{frame}




\section{Virtuální paměť}


\begin{frame}
\frametitle{Virtualizace paměti}

Cíle virtualizace:
\begin{itemize}
 \item Ochrana jádra operačního systému
 \begin{itemize}
 \item část virtuálního prostoru je chráněna proti přístupu v U - módu
 \item tato část prostoru je přístupná pouze z S nebo M módu
\end{itemize}
\item Izolace procesů
 \begin{itemize}
 \item Každý proces - běžící program - má svůj virtuální prostor.
 \item Proces se pohybuje pouze ve virtuálním prostoru.
 \item Součástí virtuálního prostoru je i paměť rezervovaná pro OS, aby mohl procesor skočit na obsluhu vyjímky a přerušení.
\end{itemize}
\end{itemize}
\end{frame}


\begin{frame}
\frametitle{Virtualizace paměti}

Další cíle virtualizace:
\begin{itemize}
 \item Efektivní využití RAM
 \begin{itemize}
 \item virtuální prostory všech procesů je potřeba poskládat do fyzické paměti - RAM
 \item ukončením procesu vznikají díry volné paměti RAM různých velikostí.
\end{itemize}
\item Podpora sdílení paměti
 \begin{itemize}
 \item Text programů (strojové instrukce) mohou být sdílené různými procesy, které vykonávají stejný program.
 \item Text knihoven (strojové instrukce) jsou sdíleny, aby se ušetřila fyzická paměť.
 \end{itemize}
\item Mapování zařízení a souborů do virtuálního procesu
 \begin{itemize}
 \item Pro rychlejší přístup k zařízení je možné fyzickou adresu periferie mapovat do virtuálního prostoru procesu.
 \item Při čtení nebo zápisu do souboru se obsah kopíruje do paměti OS, pro zrychlení je možné mapovat tuto část OS do virtuálního prostoru procesu.
 \end{itemize}
\end{itemize}
\end{frame}


\begin{frame}
\frametitle{Virtualizace paměti}

Základní princip stránkování jako nejpoužívanější způsob virtualizace paměti
\begin{itemize}
\item paměť se rozdělí na stejně veliké části
\begin{itemize}
\item část virtuální paměti se nazývá stránka - page.
\item část fyzické paměti se nazývá rámec - frame.
\end{itemize}
\item velikost stránky a rámce je stejná a bývá 4 kiB tedy 4096 bajtů.
\item stránku lze umístit do libovolného rámce
\item umístění stránky do rámce se definuje pomocí tabulky stránek (page table):
\begin{itemize}
 \item tabulka má položku pro každou stránku
 \item položka obsahuje:
\begin{itemize}
 \item příznaky zda je stránka ve fyzické paměti
 \item příznaky přístupu R, W, X, U
 \item číslo rámce fyzické paměti
\end{itemize}
\end{itemize}
\end{itemize}
\end{frame}


\begin{frame}
\frametitle{Virtualizace paměti}

Princip stránkování
\begin{itemize}
\item paměť se rozdělí od počátku na stejně velké části, většinou $4kiB = 4096B$
\begin{itemize}
\item každá část paměti začíná na adrese $k \cdot 4096$ a končí na adrese $k \cdot 4096 + 4095$
\item pro zjištění čísla stránky se využije operace $adr // 4096 = adr >> 12$
\end{itemize}
\item velikost tabulky je dána velikostí virtuálního prostoru
\begin{itemize}
\item velikost 32-bitového virtuálního prostoru je $2^{32} = 4\text{GiB}$
\item počet stránek $\frac{2^{32}}{2^{12}} = 2^{20} = 1\text{Mi} = 1024 \cdot 1024$
\item velikost jedné položky v tabulce jsou 4 bajty
\item velikost celé tabulky stránek je 4 MiB
\end{itemize}
\item velikost tabulky v 64-bitovém prostoru
\begin{itemize}
\item velikost 64-bitového virtuálního prostoru je $2^{64} = 16\text{EiB}$ (16 Exbibajtů = 16 * 1024 Pebibajtů)
\item počet stránek $\frac{2^{64}}{2^{12}} = 2^{52} = 4\text{Pi}$
\item velikost jedné položky je 8 bajtů
\item velikost tabulky stránek je 32 PiB
\end{itemize}
\end{itemize}
\end{frame}



\begin{frame}
\frametitle{Virtualizace paměti}

Stránkovací tabulka
\begin{itemize}
\item velikost stránkovací tabulky je nepřiměřeně velká
\item většina položek je zbytečná, obsahuje samé 0
\begin{itemize}
\item pouze stránky, které jsou použity pro program, nebo data jsou mapovány do fyzické paměti
\item ostatní stránky nejsou použity a nejsou mapovány do fyzické paměti
\end{itemize}
\item tabulka stránek se rozdělí na stejně veliké podtabulky a přidá se nová úroveň -- tabulka, která bude určovat, zda a kde odpovídající podtabulka existuje
\item Adresní prostor Sv32 -- 32bitové ukazatele, 32bitový prostor
\begin{itemize}
\item počet stránek $\frac{2^{32}}{2^{12}} = 2^{20} = 1\text{Mi} = 1024 \cdot 1024$
\item rozdělíme tabulku stránek na 1024 malých tabulek, každá obsahuje 1024 položek
\item pokud by malá tabulka měla pouze nulové hodnoty, nemusí existovat a zabírat paměť.
\item přidáme další nadřazenou tabulku, která obsahuje 1024 položek, zda existuje odpovídající malá tabulka a ve kterém rámci tato malá tabulka je
\begin{itemize}
\item všiměte si, že velikost malé tabulky i velikost nadřazené tabulky je přesně velikost rámce!
\end{itemize}
\end{itemize}
\end{itemize}
\end{frame}


\begin{frame}
\frametitle{Virtualizace paměti}

Položka stránkovací tabulky
\begin{itemize}
\item položka má velikost 4 bajty tj. 32 bitů
\item potřebujeme zakódovat číslo rámce:
\begin{itemize}
\item číslo rámce má v Sv32 -- 22 bitů
\item to znamená, že velikost fyzické paměti může být až $12 + 22 = 34$ tedy 16 GiB
\item K čemu je to dobré, když program může využít maximálně 4 GiB (každá virtuání adresa má pouze 32 bitů)?
\begin{itemize}
\item V počítači může běžet hodně programů a celkové využití paměti tedy lehce přesáhne 4 GiB
\end{itemize}
\end{itemize}
\item také potřebujeme zakódovat přístupová práva ke stránce -- příznaky:
\begin{itemize}
\item Příznak V - validity, nejdůležitější příznak definující, že číslo rámce je validní
\item Příznaky R,W,X,U - práva pro čtení, zápis, spuštění instrukce, přístup pro uživatele v U módu
\item Další příznaky slouží k zjištění používání stránky - A accessed, D dirty.
\end{itemize}
\end{itemize}
\end{frame}



\begin{frame}
\frametitle{Virtualizace paměti - příklad}

Nejdříve musíme vytvořit a inicializovat stránkovací tabulku
\begin{itemize}
\item Vytvoříme si dvě místa v paměti, která musí začínat na násobcích 4096 a obsahovat 4096 bajtů
\item První bude pro tabulku úrovně 1, tedy nadřazenou tabulku
\item Druhá bude pro tabulku úrovně 0, tabulku, která už obsahuje informace o mapování virtuální paměti do fyzické paměti
\end{itemize}

Kolik paměti maximálně můžeme namapovat pomocí těchto dvou tabulek
\begin{itemize}
 \item[A] celý virtuální prostor
 \item[B] 4 MiB
 \item[C] 1 MiB
 \item[D] 4 kiB
\end{itemize}
\end{frame}


\begin{frame}
\frametitle{Virtualizace paměti - příklad}

Nejdříve musíme vytvořit a inicializovat stránkovací tabulku
\begin{itemize}
\item Vytvoříme si dvě místa v paměti, která musí začínat na násobcích 4096 a obsahovat 4096 bajtů
\item První bude pro tabulku úrovně 1, tedy nadřazenou tabulku
\item Druhá bude pro tabulku úrovně 0, tabulku, která už obsahuje informace o mapování virtuální paměti do fyzické paměti
\begin{itemize}
\item Tabulky propojíme tak, že do tabulky úrovně 1 na první 4 bajty zapíšeme pozici tabulky úrovně 0.
\item VPN = (adr(tab\_0) > > 12) < < 10
\item a nastavení příznaku V, bity R,W,X musí být nulové (použito pro velké stránky 4 MiB).
\end{itemize}
\end{itemize}

Kterou část virtuální paměti můžeme díky tomuto propojení využít
\begin{itemize}
 \item[A] celý virtuální prostor
 \item[B] 0 -- 4 MiB
 \item[C] 4 MiB -- 8 MiB
 \item[D] 8 MiB -- 12 MiB
\end{itemize}
\end{frame}


\begin{frame}
\frametitle{Virtualizace paměti - příklad}

Nyní si do virtuálního prostoru namapujeme na stejné adresy program:
\begin{itemize}
\item Protože program začíná na adrese 0x200 a vejde se do 3584 bajtů (0x1000-0x200=0xE00=3584), tak stačí namapovat stránku 0 do rámce 0
\item VPN = 0
\item práva musíme nastavit X pro možnost vykonávání programu
\end{itemize}

Náš ukázkový program kopíruje data z virtuální adresy 0x6000 na adresu 0x7000. Fyzicky jsou ale zdrojová data na adrese 0x1000 a kopírovat se budou do fyzické paměti 0x2000.
\begin{itemize}
\item Pro data si stránku 6 namapujeme do rámce 1 a stránku 7 do rámce 2
\item práva můžeme nastavit pro data R,W pro obě stránky
\begin{itemize}
\item Pokud by zdroj dat byly konstanty, můžeme pro zdrojovou stránku nastavit pouze práva pro čtení a tím ochráníme data proti zápisu
\end{itemize}
\end{itemize}

\end{frame}


\begin{frame}
\frametitle{Virtualizace paměti - příklad}

Posledním krokem je nastavení módu stránkování Sv32:
\begin{itemize}
\item v registru satp (Supervisor Address Translation and Protection) nastavíme mód na vybranou hodnotu, tedy Sv32
\item jakmile je registr satp nastaven na nenulovou hodnotu módu, překlad adres a ochrana paměti je spuštěna pro S a U mód (M mód má neustále přímý přístup do paměti)
\end{itemize}

Po přepnutí do U módu spustíme náš program, který bude mít text programu na virtuální adrese 0x200 a data na virtuálních adresách 0x6000 a 0x7000. Žádné další stránky nejsou namapovány, takže nelze žádný jiný program, než náš spustit
\begin{itemize}
\item Přepnutí do U/S módu provedeme instrukcí mret tak, že nejdříve nastavíme na jakou adresu se přepneme
\item dále nastavíme, který mód po mret bude vybrán se přepneme
\end{itemize}

Nyní již náš program běží s virtualizací paměti, ve které vidí pouze program a zdrojová a cílová data.
\end{frame}



\begin{frame}
\frametitle{Virtualizace paměti - příklad}

Jeden rámec fyzické paměti může být na více místech virtuální paměti
\begin{itemize}
\item V navazujícím předmětu OSY si ukážeme praktické použití
\begin{itemize}
\item Například při alokaci velké paměti a její inicializaci na 0 může OS v rámci šetření alokovat a inicializovat pouze jeden rámec na 0 a ve virtuálním prostoru nastaví všechny alokované stránky na tento jeden rámec
\item Rámec se musí nastavit jako read only, tedy bez příznaku W, a zápis do této virtuální paměti vytvoří chybu na jejímž základě OS zkopíruje data do nového rámce a pro tento danný přístup již nastaví ve virtuálním prostoru pouze pro tuto stránku nový rámec
\item Díky tomu se zátěž alokace a inicializace rozloží na dobu, kdy se do alokované paměti bude zapisovat. Pro čtení všichni vidí ty samé 0.
\end{itemize}
\item My si ukážeme, že cílovou paměť lze namapovat ještě na stránku např. 0x8000.
\item Při kopírování dat do cílové stránky se automaticky modifikuje i stránka 0x8000, protože sdílejí stejnou fyzickou paměť, která jediná obsahuje informace.
\end{itemize}

\end{frame}




\begin{frame}
\frametitle{Efektivita tabulky stránek}
\begin{itemize}
\item každý přístup do paměti znamená, že je potřeba převést číslo stránky na číslo rámce, to znamená číst z RAM
\item jedna instrukce může číst i dvakrát z paměti z různých stránek
\item jedna instrukce tedy potřebuje 4 přístupy do paměti
\item přístup do pamět velmi zdržuje
\end{itemize}

\begin{center}
\includegraphics[width=0.7\textwidth]{generated/memory/paging/no-tbl.pdf}
\end{center}
\end{frame}

\begin{frame}
\frametitle{TLB}
\begin{itemize}
\item Řešení -- speciální rychlá cache paměť pro čísla rámců a čísel stránek -- Translation Lookaside Buffer
\item asociativní paměť -- paměť adresovaná obsahem
\begin{itemize}
\item oproti normální paměti, ptám se, kde v paměti je hodnota 15?
\end{itemize}
\item TLB je asociativní paměť
\item relativně malá kapacita, vysoká efektivita a zrychlení přístupu do paměti
\item nevýhoda -- obvodová složitost implementace TLB
\item MMU se zeptá TLB znáš hodnotu rámce pro číslo stránky p? Odpověď buď ano je to $r$ (TLB hit), nebo neznám (TLB miss)
\end{itemize}
\end{frame}



\begin{frame}
\frametitle{Tabulka stránek}
\begin{center}
\includegraphics[width=0.9\textwidth]{generated/memory/paging/page-table.pdf}
\end{center}
\end{frame}




\begin{frame}
\frametitle{Význam TLB}
\begin{itemize}
\item Skutečná přístupová doba -- Effective Access Time (EAT) -- 10--100 cyklů procesoru
‏\item Přístupová doba TLB -- 0.5 -- 1 cyklus procesoru
\item Neúspěšnost TLB, TLB miss -- 0.01 \% -- 1 \%
\item Příklad:
\begin{itemize}
\item Přístupová doba do fyzické paměti $t=30 cyklu$, do TLB $t_{TLB} = 1 cyklus$
\item Neúspěšnost TLB $\alpha = 0.01$ (jedno procento)
\item Stránkování bez TLB $ t_{celkem} = 2\cdot t = 60 cyklu$
\item Průměrná doba přístupu do paměti s TLB $ t_{celkem} = \alpha \cdot (t_{TLB}+2\cdot t) + (1-\alpha)\cdot (t_{TLB}+t) = 31.3 cyklu $
\end{itemize}
\end{itemize}
\end{frame}



\begin{frame}
\frametitle{Vliv TLB}
\begin{itemize}
\item Velikost TLB 8--4096 položek
\item Moderní procesory mají více-úrovní TLB – podobně jako úrovně cache
\item Intel Core i7 má 64 TLB položek L1 první úrovně a 1536 TLB položek L2 úrovně
\item I pro malé TLB je úspěšnost nalezení položky 99--99.99\% (souvisí s principem lokality tj. prostorovou závislostí programů)
\item Problém TLB je při změně procesu a tím i změně tabulky stránek
\item Intel umožňuje speciálními instrukcemi ponechat stránku v TLB
\end{itemize}
\end{frame}



\begin{frame}
\frametitle{64-bitový mód}

Co je nutné změnit pro 64-bitový mód?
\begin{itemize}
\item Rozsah čísla rámce
\begin{itemize}
\item V Sv32 je k dispozici 22 bitů pro číslo rámce, což umožňuje využít 16 GiB fyzické paměti (to není mnoho)
\item Pro 64-bitový mód potřebujeme více bitů na číslo rámce, ale potřebujeme zachovat minimálně 10 bitů práv
\begin{itemize}
\item Řešení: Velikost jedné položky v tabulce stránek bude 64 bitů tj. 8 bajtů
\end{itemize}
\end{itemize}
\item Velikost tabulky
\begin{itemize}
\item V Sv32 je velikost virtuálního prostoru $2^{32}$ tj. 4 GiB, počet položek tabulky stránek je $2^{20} = 1024 \cdot 1024$ a velikost tabulky je 4 MiB
\item Pro 64-bitový mód je velikost virtuálního prostoru $2^{64}$, počet položek tabulky stránek je $2^{52}$ a velikost tabulky je 32 PiB
\begin{itemize}
\item Řešení: Omezíme velikost prostoru a zavedeme více úrovní stránkovacích tabulek
\end{itemize}
\end{itemize}
\end{itemize}
\end{frame}


\begin{frame}
\frametitle{Stránkování Sv39}

\begin{itemize}
\item velikost stránky je 4 KiB tj. $2^{12}$
\item každá tabulka všech úrovní se vejde přesně do jedné stránky
\begin{itemize}
\item Velikost jedné položky v tabulce stránek je 8 bajtů, tzn. celkem se do jedné stránky vejde $512=2^{9}$ položek
\end{itemize}
\item Stránkování Sv39 bude používat pouze spodních 39 bitů z adresy, tj. velikost virtuálního prostoru je 512 GiB.
\begin{itemize}
\item Existují 3-úrovně stránkovacích tabulek - kódováno $3 \cdot 9$ bity
\item 12-bitů kóduje polohu uvnitř stránky
\item Celkem tedy 39 bitů virtuální adresy
\end{itemize}
\end{itemize}

Proč nevzít rovnou 64 bitů virtuálního prostoru?
\end{frame}


\begin{frame}
\frametitle{Stránkování Sv48, Sv 57}

\begin{itemize}
\item Každá další úroveň přináši zpomalení, protože je nutné načítat další informaci z paměti.
\item Stránkování Sv48 zavádí 4-úrovně stránkovacích tabulek
\begin{itemize}
\item Kódováno $4 \cdot 9 + 12 = 48$ bitů adresy
\item Velikost virtuálního prostoru je 256 TiB.
\end{itemize}
\item Stránkování Sv57 zavádí 5-úrovní stránkovacích tabulek
\begin{itemize}
\item Kódováno $5 \cdot 9 + 12 = 57$ bitů adresy
\item Velikost virtuálního prostoru je 128 PiB.
\end{itemize}
\end{itemize}

\end{frame}



\end{document}
